\documentclass[11pt]{article}
\usepackage[utf8]{inputenc}
\usepackage{amsmath}
\usepackage{amssymb}
\usepackage{geometry}
\usepackage{booktabs}
\usepackage{hyperref}
\geometry{margin=1in}

\hypersetup{colorlinks=true, linkcolor=blue, urlcolor=blue}

\title{MethylUtils Theoretical Foundation}
\author{MethylPipeline}
\date{}

\begin{document}
\maketitle
\tableofcontents
\newpage

% -----------------------------------------------------------------------
\section{Role}

\texttt{MethylUtils} is the shared mathematical layer used by all other
packages in MethylPipeline for centroid comparison, ECDF overlap, significance
testing, and the canonical biological \texttt{effect\_size}.

% -----------------------------------------------------------------------
\section{Per-Position Notation}

For two centroid groups (e.g.\ healthy and disease), let the per-position
quantities at genomic position $i$ be

\begin{itemize}
  \item \(N_1,\,N_2\) --- number of contributing samples in each group.
  \item \(\hat\mu_1,\,\hat\mu_2 = S_{x,k}/N_k\) --- mean methylation fractions.
  \item \(\hat\sigma_1^2,\,\hat\sigma_2^2\) --- sample variances
        (Equation~\eqref{eq:samplevar}).
\end{itemize}

Sample variances are derived from stored sufficient statistics without
reading individual sample files:

\begin{equation}
  \hat\sigma_k^2 = \frac{S_{x^2,k} - S_{x,k}^2/N_k}{\max(N_k-1,\,1)}.
  \label{eq:samplevar}
\end{equation}

% -----------------------------------------------------------------------
\section{Mann--Whitney $U$ from Bin Counts}

The core significance test is \texttt{mann\_whitney\_from\_bin\_counts()},
which approximates the Mann--Whitney $U$ statistic directly from the centroid
\texttt{bin\_counts} histograms, avoiding the need for raw sample values.

Let \(\mathrm{bc}_{k,b}\) be the count in bin $b$ for group $k$, and
\(\mathrm{CS}_{2,b} = \sum_{b'<b}\mathrm{bc}_{2,b'}\) the cumulative count of
group 2 below bin $b$.  Then

\begin{equation}
  U = \sum_{b} \mathrm{bc}_{1,b}\,\mathrm{CS}_{2,b}
      + \frac{1}{2}\sum_{b} \mathrm{bc}_{1,b}\,\mathrm{bc}_{2,b}.
\end{equation}

The z-statistic corrects for ties using \(t_b = \mathrm{bc}_{1,b}+\mathrm{bc}_{2,b}\):

\begin{equation}
  \mathrm{Var}(U) = \frac{N_1 N_2}{12}
    \left[(N_1+N_2+1) - \frac{\sum_b t_b(t_b^2-1)}{(N_1+N_2)(N_1+N_2-1)}\right],
  \qquad z = \frac{U - N_1 N_2/2}{\sqrt{\mathrm{Var}(U)}}.
\end{equation}

% -----------------------------------------------------------------------
\section{Between-Sample Heterogeneity}

\texttt{dl\_heterogeneity()} provides a DerSimonian--Laird-style estimate of
between-sample variance using only the stored centroid statistics
\((\Sigma m,\,\Sigma u,\,\Sigma w x^2,\,\Sigma c^2,\,N)\).

Let \(C = \Sigma m + \Sigma u\).  The between-sample variance is

\begin{equation}
  \hat\tau^2 = \max\!\left(0,\;
    \frac{\displaystyle\Sigma w x^2 - \frac{(\Sigma m)^2}{C}
          - (N-1)}{C - \Sigma c^2/C}
  \right).
\end{equation}

Large \(\hat\tau^2\) suggests heterogeneous sub-populations rather than a
uniform methylation shift; MethylDetector can use this to filter noisy loci.

% -----------------------------------------------------------------------
\section{Continuous ECDF Overlap}

Centroids store per-position histograms (\texttt{bin\_counts}).  The
\texttt{ECDFView} class reconstructs a continuous CDF over \([0,1]\) using
PCHIP interpolation and derives a PDF by differentiating the spline.

\paragraph{Definition.}  Overlap between two class distributions is

\begin{equation}
  O = \int_0^1 \min\!\bigl(f_1(x),\,f_2(x)\bigr)\,dx,
  \label{eq:overlap}
\end{equation}

where \(f_k(x) = F_k'(x)\) is the PCHIP-derived PDF for class $k$.

The integral is evaluated with the trapezoidal rule on a grid of 512 points.
PDFs are renormalised before integration:

\begin{equation}
  \tilde{f}_k(x) = \frac{f_k(x)}{\displaystyle\int_0^1 f_k(t)\,dt}.
\end{equation}

\(O \in [0,1]\): \(O=1\) means identical distributions; \(O=0\) means no
shared support.

Implemented in \texttt{ecdf\_overlap\_integral()}.

\paragraph{Discrete Bhattacharyya approximation.}

For large position sets, a fast approximation uses normalised bin counts:

\begin{equation}
  O_\mathrm{Bh} = \sum_{b=1}^{B} \sqrt{p_{1b}\cdot p_{2b}},
  \qquad p_{kb} = \frac{\mathrm{bc}_{kb}}{\sum_{b'}\mathrm{bc}_{kb'}}.
\end{equation}

Implemented in \texttt{discrete\_overlap\_from\_bin\_counts()}.

% -----------------------------------------------------------------------
\section{Variance Penalty (Reliability Term)}

The reliability term in the canonical effect size is

\begin{equation}
  R = \exp\!\bigl(-\lambda_\mathrm{var}
            (\sqrt{\hat\sigma_1^2} + \sqrt{\hat\sigma_2^2})\bigr).
  \label{eq:reliability}
\end{equation}

\begin{itemize}
  \item \(\lambda_\mathrm{var}\) (default \(2.0\)) controls the penalty
        strength.
  \item The two group variances are treated independently; no equal-variance
        assumption is made.
  \item A position with tight within-group distributions has \(R\approx 1\);
        a diffuse position has \(R \ll 1\).
  \item These variances affect only the biological reliability term; the
        significance gate itself remains non-parametric.
\end{itemize}

% -----------------------------------------------------------------------
\section{Canonical Biological Effect Size}

The single score used for DMP ranking and selection is

\begin{equation}
  \boxed{
    \mathrm{effect\_size}
    = |\Delta\mu| \cdot (1 - O) \cdot R,
  }
  \label{eq:effectsize}
\end{equation}

where \(|\Delta\mu| = |\hat\mu_1-\hat\mu_2|\), \(O\) is the ECDF overlap
(Equation~\eqref{eq:overlap}), and \(R\) is the reliability term
(Equation~\eqref{eq:reliability}).

The score is clipped to \([0,1]\).  Implemented in
\texttt{effect\_size\_from\_components()} and the convenience wrapper
\texttt{ecdf\_effect\_size()}.

% -----------------------------------------------------------------------
\section{Downstream Consistency}

\texttt{MethylDetector} and \texttt{MethylDetectorExplorer} both delegate
their final overlap and effect-size computation to these shared MethylUtils
helpers.  This ensures that the score used to \emph{select} DMPs is the same
formula used to \emph{rank} positions in exploratory analysis.

\end{document}
